\section{课程背景与核心问题}

本讲主题是 \textbf{Post-Training Verifiable Agents}。Jiantao Jiao 的切入角度非常明确：如果我们希望 Agent 在真实环境里长期稳定执行任务，仅靠 ``对话质量好'' 不够，必须把训练目标从 ``讨好人类偏好'' 扩展为 ``在环境反馈下可验证地完成任务''。也就是说，模型输出不再只是文本，而是带状态影响的行动轨迹（trajectory）。

\begin{importantbox}{本讲主命题}
Agentic 模型的后训练目标可以写成一句话：\textbf{在保持人类可用性的前提下，最大化可验证任务回报（verifiable rewards）}。它要求系统同时满足 ``会沟通'' 与 ``会做事'' 两类能力，而不是只优化其中一侧。
\end{importantbox}

讲者在开场就强调，很多早期聊天模型主要优化的是人类偏好（preference alignment），例如回答礼貌、语气自然、内容看起来合理。但 Agent 任务往往包含代码修改、工具调用、数据库查询、外部 API 交互等动作，这些动作会改变环境状态，错误也会有累计效应。因此，训练目标需要显式纳入环境反馈。

\begin{knowledgebox}{从 Chatbot 到 Agent 的目标函数变化}
\begin{itemize}
    \item Chatbot：偏向 ``回答是否让人满意''，奖励常来自偏好比较。
    \item Agent：偏向 ``任务是否完成且可验证''，奖励更多来自程序化 verifier。
    \item 两者并非对立：生产系统通常需要 ``偏好质量 + 任务正确率'' 的联合优化。
\end{itemize}
\end{knowledgebox}

课程还反复指出一个实践问题：当模型开始执行复杂任务时，``看起来像对'' 与 ``真的做对'' 差距会迅速扩大。越是多步任务，这个差距越明显。因此，是否有高质量 verifier，以及 verifier 是否覆盖关键失败模式，直接决定后训练上限。

\subsection{本章小结}
本章建立了后续全部讨论的前提：Agent 后训练不是把聊天模型 ``再微调一下''，而是把训练目标重构为 ``环境中的可验证任务完成''。这会牵引数据、评估和算法三条线同时变化。

